\vspace{0.5cm}
\section{Applicazione software di accesso al database}

Per sfruttare efficacemente l'architettura del database, è stato
sviluppato un client da riga di comando (CLI), eseguibile compilando
il file \texttt{client.c}, scritto in C.

Per aumentare la portabilità e l'indipendenza dall'ambiente, al lancio
l'applicativo richiede l'immissione delle credenziali di accesso al 
database PostgreSQL. Una volta verificata e stabilita la connessione con
il server, il programma permette l'interrogazione del database,
scegliendo tra le seguenti query:
\begin{enumerate}
    \item Individuazione del modello di veicolo più venduto.
    \item Estrazione dei venditori che hanno superato una specifica
          soglia di vendite in un determinato intervallo temporale.
    \item Ricerca del ricambio più utilizzato per ogni tipologia
          di motore, con ordinamento decrescente in base alla frequenza d'uso.
    \item Calcolo del numero di vendite per le combinazioni più richieste di
          frazionamento e alimentazione.
    \item Calcolo del prezzo medio dei veicoli venduti soggetti alla
          tassa del superbollo (sopra i 250 cv).
\end{enumerate}

La selezione della query desiderata avviene tramite l'inserimento del codice numerico corrispondente. 

La seconda interrogazione è stata implementata in forma parametrica: a tempo di esecuzione,
il sistema richiede all'utente di inserire il valore soglia relativo al volume di vendite.
Questa soluzione si rivela particolarmente funzionale in scenari gestionali, come la definizione
di politiche incentivanti o l'assegnazione di bonus basati sulle performance, consentendo di adattare
dinamicamente il target richiesto in risposta alle fluttuazioni stagionali del mercato
(ad esempio, innalzando gli obiettivi nei periodi di maggiore affluenza).